\documentclass[11pt,a4paper]{article}
\usepackage{fontspec,amsmath,amssymb,amsthm,microtype}
\setmainfont{lmroman10-regular.otf}[BoldFont=lmroman10-bold.otf,ItalicFont=lmroman10-italic.otf,BoldItalicFont=lmroman10-bolditalic.otf]
\usepackage{xcolor}
\usepackage[most]{tcolorbox}
\usepackage{mathrsfs}
\providecommand{\lt}{<}
\providecommand{\gt}{>}
\setmonofont{DejaVuSansMono.ttf}[BoldFont=DejaVuSansMono-Bold.ttf,ItalicFont=DejaVuSansMono-Oblique.ttf,BoldItalicFont=DejaVuSansMono-BoldOblique.ttf,Scale=0.86]
\newfontfamily\leanfallback{STIXTwoMath-Regular.otf}[Scale=0.86]
\definecolor{leanrule}{gray}{0.62}
\definecolor{leanback}{gray}{0.962}
\definecolor{leantitle}{gray}{0.22}
\definecolor{leanmuted}{gray}{0.38}
\newtcolorbox{leanbox}[1]{enhanced,breakable,colback=leanback,colframe=leanrule,
  boxrule=0.4pt,arc=0.6mm,left=2.4mm,right=2.4mm,top=1.8mm,bottom=1.8mm,
  before skip=0.8em,after skip=1em,fonttitle=\small\bfseries,coltitle=leantitle,
  colbacktitle=leanback,titlerule=0.3pt,title={#1}}
\newlength{\leancharwidth}
\newenvironment{leancode}{\par\smallskip\begingroup\ttfamily\footnotesize\raggedright
  \settowidth{\leancharwidth}{M}\setlength{\parskip}{0pt}\setlength{\parindent}{0pt}
  \frenchspacing}{\par\endgroup\smallskip}
\newcommand{\leanline}[2]{\hangindent=\dimexpr#1\leancharwidth+4\leancharwidth\relax
  \hangafter=1\noindent\hspace*{#1\leancharwidth}#2\par}
\newcommand{\leansource}[1]{\par{\footnotesize\color{leanmuted}#1}\par}
\newcommand{\leanlabel}[1]{\par\smallskip\noindent{\small\bfseries #1}\par}
\usepackage[margin=25mm]{geometry}
\usepackage[unicode,hidelinks]{hyperref}
\hypersetup{pdfauthor={LangevinLean Project}}
\setlength{\emergencystretch}{3em}
\setlength{\parskip}{0.4em}
\setlength{\parindent}{0pt}
\allowdisplaybreaks[2]
\raggedbottom
\title{\texorpdfstring{The linear IPLA discretization term and explicit accuracy arithmetic}{The linear IPLA discretization term and explicit accuracy arithmetic}}
\author{LangevinLean Project}
\date{}
\begin{document}\maketitle
Draft mathematical exposition; these notes may contain mistakes.


Let \(N>0\), \(0<m\le L\), and \(D=d_x+d_\theta/N\). The lifted sharp bias simplifies exactly to \(B=6\gamma\sqrt{(L/m)LD}\). After multiplication by \(\sqrt m\), it is \(6L\gamma\sqrt D\). These scalar identities yield explicit particle, step and iteration choices for normalized parameter accuracy, and an exact small-step comparison with Project 03. The probabilistic bounds and their assumptions are proved in the application modules; they are not supplied or presumed by these arithmetic lemmas.


\section{\texorpdfstring{Exact particle scaling and normalization}{Exact particle scaling and normalization}}


\textbf{Definitions and claims.} Set \(D=d_x+d_\theta/N\ge0\). The new and Project 03 discretization terms are

\[B(\gamma)=6\gamma\sqrt{\frac Lm\,LD},\qquad B_{03}(\gamma)=\frac{33}{20}\frac Lm\sqrt D\sqrt\gamma.\]


The exact lifted expression satisfies

\[6\frac\gamma N\sqrt{\frac{NL}{Nm}(NL)(d_\theta+Nd_x)}=B(\gamma).\]


Moreover, for \(m>0\), \(L,D\ge0\),

\[\sqrt m\,B(\gamma)=6L\gamma\sqrt D,\qquad \sqrt m\sqrt{\frac{2d_\theta}{Nm}}=\sqrt{\frac{2d_\theta}{N}}.\]


\textbf{Proof.} Cancel the positive particle number in the condition number. Since \(d_\theta+Nd_x=ND\), the expression under the remaining square root is \(N^2(L/m)LD\). Extracting its nonnegative factor \(N\) cancels the denominator outside. For the normalized bias, multiply the radicand by \(m\) to obtain \(L^2D\) and extract \(L\ge0\). For concentration, the same multiplication cancels \(m\) under the square root. All divisions by the particle number and curvature use their proved positivity.

\begin{leanbox}{Lean: \texttt{iplaEffectiveDimension}, \texttt{iplaLinearBias}, \texttt{iplaProject03Bias} and 4 more}
\leansource{\texttt{iplaEffectiveDimension} \quad WassersteinMixing/\allowbreak{}IPLARates.lean, lines 17–18}
\begin{leancode}
\leanline{0}{def\ iplaEffectiveDimension\ (N\ dθ\ dx\ :\ ℕ)\ :\ ℝ\ :=}
\leanline{2}{(dx\ :\ ℝ)\ +\ (dθ\ :\ ℝ)\ /\ N}
\end{leancode}
\leansource{\texttt{iplaLinearBias} \quad WassersteinMixing/\allowbreak{}IPLARates.lean, lines 21–22}
\begin{leancode}
\leanline{0}{def\ iplaLinearBias\ (m\ L\ γ\ D\ :\ ℝ)\ :\ ℝ\ :=}
\leanline{2}{6\ *\ γ\ *\ Real.sqrt\ ((L\ /\ m)\ *\ L\ *\ D)}
\end{leancode}
\leansource{\texttt{iplaProject03Bias} \quad WassersteinMixing/\allowbreak{}IPLARates.lean, lines 25–26}
\begin{leancode}
\leanline{0}{def\ iplaProject03Bias\ (m\ L\ γ\ D\ :\ ℝ)\ :\ ℝ\ :=}
\leanline{2}{(33\ /\ 20\ :\ ℝ)\ *\ (L\ /\ m)\ *\ Real.sqrt\ D\ *\ Real.sqrt\ γ}
\end{leancode}
\leansource{\texttt{iplaEffectiveDimension\_\allowbreak{}nonneg} \quad WassersteinMixing/\allowbreak{}IPLARates.lean, lines 44–47}
\begin{leancode}
\leanline{0}{theorem\ iplaEffectiveDimension\_nonneg\ (N\ dθ\ dx\ :\ ℕ)\ :}
\leanline{4}{0\ ≤\ iplaEffectiveDimension\ N\ dθ\ dx}
\end{leancode}
\leansource{\texttt{ipla\_\allowbreak{}lifted\_\allowbreak{}bias\_\allowbreak{}eq} \quad WassersteinMixing/\allowbreak{}IPLARates.lean, lines 50–65}
\begin{leancode}
\leanline{0}{theorem\ ipla\_lifted\_bias\_eq\ (N\ dθ\ dx\ :\ ℕ)\ (hN\ :\ 0\ <\ N)\ (m\ L\ γ\ :\ ℝ)\ :}
\leanline{4}{6\ *\ (γ\ /\ (N\ :\ ℝ))\ *}
\leanline{6}{Real.sqrt\ ((((N\ :\ ℝ)\ *\ L)\ /\ ((N\ :\ ℝ)\ *\ m))\ *}
\leanline{8}{((N\ :\ ℝ)\ *\ L)\ *\ (dθ\ +\ N\ *\ dx\ :\ ℕ))\ =}
\leanline{6}{iplaLinearBias\ m\ L\ γ\ (iplaEffectiveDimension\ N\ dθ\ dx)}
\end{leancode}
\leansource{\texttt{ipla\_\allowbreak{}linear\_\allowbreak{}bias\_\allowbreak{}scale} \quad WassersteinMixing/\allowbreak{}IPLARates.lean, lines 68–75}
\begin{leancode}
\leanline{0}{theorem\ ipla\_linear\_bias\_scale\ (m\ L\ γ\ D\ :\ ℝ)\ (hm\ :\ 0\ <\ m)}
\leanline{4}{(hL\ :\ 0\ ≤\ L)\ (hD\ :\ 0\ ≤\ D)\ :}
\leanline{4}{Real.sqrt\ m\ *\ iplaLinearBias\ m\ L\ γ\ D\ =\ 6\ *\ L\ *\ γ\ *\ Real.sqrt\ D}
\end{leancode}
\leansource{\texttt{ipla\_\allowbreak{}concentration\_\allowbreak{}scale} \quad WassersteinMixing/\allowbreak{}IPLARates.lean, lines 78–84}
\begin{leancode}
\leanline{0}{theorem\ ipla\_concentration\_scale\ (N\ dθ\ :\ ℕ)\ (hN\ :\ 0\ <\ N)\ (m\ :\ ℝ)\ (hm\ :\ 0\ <\ m)\ :}
\leanline{4}{Real.sqrt\ m\ *\ Real.sqrt\ (2\ *\ (dθ\ :\ ℝ)\ /\ ((N\ :\ ℝ)\ *\ m))\ =}
\leanline{6}{Real.sqrt\ (2\ *\ (dθ\ :\ ℝ)\ /\ N)}
\end{leancode}
\leanlabel{How the Lean proof follows the argument}
The first three definitions retain the literal dimension and the two constants. \texttt{ipla\allowbreak{}Effective\allowbreak{}Dimension\_\allowbreak{}nonneg} proves the dimension sign. \texttt{ipla\_\allowbreak{}lifted\_\allowbreak{}bias\_\allowbreak{}eq} factors the radicand as the square of the particle number times the reduced radicand, then uses the square-root product identity. \texttt{ipla\_\allowbreak{}linear\_\allowbreak{}bias\_\allowbreak{}scale} invokes the proved curvature scaling identity. \texttt{ipla\_\allowbreak{}concentration\_\allowbreak{}scale} combines positive square roots before exact cancellation.
\end{leanbox}

\section{\texorpdfstring{A positive step with an exact reciprocal iteration factor}{A positive step with an exact reciprocal iteration factor}}


\textbf{Definitions.} For \(L>0\) and a desired normalized accuracy \(\varepsilon>0\), set

\[M=\max\!\left\{1,\frac{18\sqrt D}{\varepsilon}\right\},\qquad \gamma_\varepsilon=\frac1{LM},\qquad T_\varepsilon=\frac Lm\,M.\]


\textbf{Claim.} The step is positive, \(\gamma_\varepsilon L\le1\), and \((m\gamma_\varepsilon)T_\varepsilon=1\).




\textbf{Proof.} The maximum is at least one and hence positive, even when \(D=0\). Thus both denominator factors are positive. Direct cancellation gives \(\gamma_\varepsilon L=1/M\le1\). Multiplying the step by its displayed iteration factor cancels the positive constants \(m,L,M\) exactly.

\begin{leanbox}{Lean: \texttt{iplaAccuracyStep}, \texttt{iplaAccuracyPrefactor}, \texttt{iplaAccuracyStep\_\allowbreak{}pos} and 2 more}
\leansource{\texttt{iplaAccuracyStep} \quad WassersteinMixing/\allowbreak{}IPLARates.lean, lines 29–30}
\begin{leancode}
\leanline{0}{def\ iplaAccuracyStep\ (L\ D\ ε\ :\ ℝ)\ :\ ℝ\ :=}
\leanline{2}{1\ /\ (L\ *\ max\ 1\ (18\ *\ Real.sqrt\ D\ /\ ε))}
\end{leancode}
\leansource{\texttt{iplaAccuracyPrefactor} \quad WassersteinMixing/\allowbreak{}IPLARates.lean, lines 33–34}
\begin{leancode}
\leanline{0}{def\ iplaAccuracyPrefactor\ (m\ L\ D\ ε\ :\ ℝ)\ :\ ℝ\ :=}
\leanline{2}{(L\ /\ m)\ *\ max\ 1\ (18\ *\ Real.sqrt\ D\ /\ ε)}
\end{leancode}
\leansource{\texttt{iplaAccuracyStep\_\allowbreak{}pos} \quad WassersteinMixing/\allowbreak{}IPLARates.lean, lines 86–89}
\begin{leancode}
\leanline{0}{theorem\ iplaAccuracyStep\_pos\ (L\ D\ ε\ :\ ℝ)\ (hL\ :\ 0\ <\ L)\ :}
\leanline{4}{0\ <\ iplaAccuracyStep\ L\ D\ ε}
\end{leancode}
\leansource{\texttt{iplaAccuracyStep\_\allowbreak{}mul\_\allowbreak{}L\_\allowbreak{}le} \quad WassersteinMixing/\allowbreak{}IPLARates.lean, lines 91–99}
\begin{leancode}
\leanline{0}{theorem\ iplaAccuracyStep\_mul\_L\_le\ (L\ D\ ε\ :\ ℝ)\ (hL\ :\ 0\ <\ L)\ :}
\leanline{4}{iplaAccuracyStep\ L\ D\ ε\ *\ L\ ≤\ 1}
\end{leancode}
\leansource{\texttt{ipla\_\allowbreak{}step\_\allowbreak{}mul\_\allowbreak{}prefactor} \quad WassersteinMixing/\allowbreak{}IPLARates.lean, lines 102–107}
\begin{leancode}
\leanline{0}{theorem\ ipla\_step\_mul\_prefactor\ (m\ L\ D\ ε\ :\ ℝ)\ (hm\ :\ 0\ <\ m)\ (hL\ :\ 0\ <\ L)\ :}
\leanline{4}{(m\ *\ iplaAccuracyStep\ L\ D\ ε)\ *\ iplaAccuracyPrefactor\ m\ L\ D\ ε\ =\ 1}
\end{leancode}
\leanlabel{How the Lean proof follows the argument}
\texttt{ipla\allowbreak{}Accuracy\allowbreak{}Step} and \texttt{ipla\allowbreak{}Accuracy\allowbreak{}Prefactor} use a maximum rather than dividing by the square root of the dimension; therefore zero dimension is included. \texttt{ipla\allowbreak{}Accuracy\allowbreak{}Step\_\allowbreak{}pos} proves positivity of the denominator. \texttt{ipla\allowbreak{}Accuracy\allowbreak{}Step\_\allowbreak{}mul\_\allowbreak{}L\_\allowbreak{}le} rewrites the scaled step as the reciprocal of that maximum. \texttt{ipla\_\allowbreak{}step\_\allowbreak{}mul\_\allowbreak{}prefactor} establishes exact equality by field cancellation using the positive factors.
\end{leanbox}

\section{\texorpdfstring{Three equal budgets give invariant-based normalized accuracy}{Three equal budgets give invariant-based normalized accuracy}}


\textbf{Claim.} Let \(R\ge0\), \(\varepsilon>0\), \(N>0\), \(0<m\le L\), and \(D\ge0\). With the preceding step, assume

\[N\varepsilon^2\ge18d_\theta,\qquad k\ge\frac Lm\max\!\left\{1,\frac{18\sqrt D}{\varepsilon}\right\}\log\!\left(1+\frac{3\sqrt m\,R}{\varepsilon}\right).\]


Then each of the following three terms is at most \(\varepsilon/3\):

\[\sqrt{\frac{2d_\theta}{N}},\qquad \sqrt m(1-m\gamma_\varepsilon)^kR,\qquad \sqrt m\,B(\gamma_\varepsilon).\]


Consequently their sum is at most \(\varepsilon\).




\textbf{Proof.} Squaring the nonnegative concentration term reduces its budget exactly to the stated particle inequality. The maximum in the step satisfies \(18\sqrt D\le\varepsilon M\), so the normalized bias \(6\sqrt D/M\) is at most \(\varepsilon/3\). Put \(t=m\gamma_\varepsilon\in(0,1]\). Its exact reciprocal factor turns the iteration threshold into \(tk\ge\log(1+3\sqrt mR/\varepsilon)\). The established geometric-transient estimate then bounds \((1-t)^k\sqrt mR\) by \(\varepsilon/3\). Adding the three inequalities gives the conclusion.

\begin{leanbox}{Lean: \texttt{ipla\_\allowbreak{}step\_\allowbreak{}bias\_\allowbreak{}budget}, \texttt{ipla\_\allowbreak{}concentration\_\allowbreak{}budget}, \texttt{ipla\_\allowbreak{}transient\_\allowbreak{}budget} and 1 more}
\leansource{\texttt{ipla\_\allowbreak{}step\_\allowbreak{}bias\_\allowbreak{}budget} \quad WassersteinMixing/\allowbreak{}IPLARates.lean, lines 110–125}
\begin{leancode}
\leanline{0}{theorem\ ipla\_step\_bias\_budget\ (m\ L\ D\ ε\ :\ ℝ)\ (hm\ :\ 0\ <\ m)\ (hL\ :\ 0\ <\ L)}
\leanline{4}{(hD\ :\ 0\ ≤\ D)\ (hε\ :\ 0\ <\ ε)\ :}
\leanline{4}{Real.sqrt\ m\ *\ iplaLinearBias\ m\ L\ (iplaAccuracyStep\ L\ D\ ε)\ D\ ≤\ ε\ /\ 3}
\end{leancode}
\leansource{\texttt{ipla\_\allowbreak{}concentration\_\allowbreak{}budget} \quad WassersteinMixing/\allowbreak{}IPLARates.lean, lines 128–135}
\begin{leancode}
\leanline{0}{theorem\ ipla\_concentration\_budget\ (N\ dθ\ :\ ℕ)\ (hN\ :\ 0\ <\ N)\ (ε\ :\ ℝ)\ (hε\ :\ 0\ <\ ε)}
\leanline{4}{(hparticles\ :\ 18\ *\ (dθ\ :\ ℝ)\ ≤\ (N\ :\ ℝ)\ *\ ε\ \textasciicircum{}\ 2)\ :}
\leanline{4}{Real.sqrt\ (2\ *\ (dθ\ :\ ℝ)\ /\ (N\ :\ ℝ))\ ≤\ ε\ /\ 3}
\end{leancode}
\leansource{\texttt{ipla\_\allowbreak{}transient\_\allowbreak{}budget} \quad WassersteinMixing/\allowbreak{}IPLARates.lean, lines 138–159}
\begin{leancode}
\leanline{0}{theorem\ ipla\_transient\_budget\ (m\ L\ D\ R\ ε\ :\ ℝ)\ (hm\ :\ 0\ <\ m)}
\leanline{4}{(hmL\ :\ m\ ≤\ L)\ (hR\ :\ 0\ ≤\ R)\ (hε\ :\ 0\ <\ ε)\ (k\ :\ ℕ)}
\leanline{4}{(hk\ :\ iplaAccuracyPrefactor\ m\ L\ D\ ε\ *}
\leanline{6}{Real.log\ (1\ +\ 3\ *\ Real.sqrt\ m\ *\ R\ /\ ε)\ ≤\ (k\ :\ ℝ))\ :}
\leanline{4}{Real.sqrt\ m\ *\ ((1\ -\ m\ *\ iplaAccuracyStep\ L\ D\ ε)\ \textasciicircum{}\ k\ *\ R)\ ≤\ ε\ /\ 3}
\end{leancode}
\leansource{\texttt{ipla\_\allowbreak{}accuracy\_\allowbreak{}arithmetic} \quad WassersteinMixing/\allowbreak{}IPLARates.lean, lines 162–173}
\begin{leancode}
\leanline{0}{theorem\ ipla\_accuracy\_arithmetic\ (N\ dθ\ :\ ℕ)\ (hN\ :\ 0\ <\ N)\ (m\ L\ D\ R\ ε\ :\ ℝ)}
\leanline{4}{(hm\ :\ 0\ <\ m)\ (hmL\ :\ m\ ≤\ L)\ (hD\ :\ 0\ ≤\ D)\ (hR\ :\ 0\ ≤\ R)\ (hε\ :\ 0\ <\ ε)}
\leanline{4}{(hparticles\ :\ 18\ *\ (dθ\ :\ ℝ)\ ≤\ (N\ :\ ℝ)\ *\ ε\ \textasciicircum{}\ 2)\ (k\ :\ ℕ)}
\leanline{4}{(hk\ :\ iplaAccuracyPrefactor\ m\ L\ D\ ε\ *}
\leanline{6}{Real.log\ (1\ +\ 3\ *\ Real.sqrt\ m\ *\ R\ /\ ε)\ ≤\ (k\ :\ ℝ))\ :}
\leanline{4}{Real.sqrt\ (2\ *\ (dθ\ :\ ℝ)\ /\ (N\ :\ ℝ))\ +}
\leanline{6}{Real.sqrt\ m\ *\ ((1\ -\ m\ *\ iplaAccuracyStep\ L\ D\ ε)\ \textasciicircum{}\ k\ *\ R\ +}
\leanline{8}{iplaLinearBias\ m\ L\ (iplaAccuracyStep\ L\ D\ ε)\ D)\ ≤\ ε}
\end{leancode}
\leanlabel{How the Lean proof follows the argument}
\texttt{ipla\_\allowbreak{}step\_\allowbreak{}bias\_\allowbreak{}budget} rewrites the normalized bias using \texttt{ipla\_\allowbreak{}linear\_\allowbreak{}bias\_\allowbreak{}scale} and the maximum inequality. \texttt{ipla\_\allowbreak{}concentration\_\allowbreak{}budget} compares squares only after proving both sides nonnegative. \texttt{ipla\_\allowbreak{}transient\_\allowbreak{}budget} applies \texttt{geometric\_\allowbreak{}transient\_\allowbreak{}le} with the reciprocal identity and the proved step range. \texttt{ipla\_\allowbreak{}accuracy\_\allowbreak{}arithmetic} adds these three derived budgets. The application later substitutes the actual nonnegative initial transport distance for the scalar radius.
\end{leanbox}

\section{\texorpdfstring{An optional conservative accuracy budget for a doubled bias}{An optional conservative accuracy budget for a doubled bias}}


\textbf{Definitions and claim.} For the target-based parameter bound, reserve room for the doubled discretization term by choosing

\[\widetilde\gamma_\varepsilon=\frac1{L\max\{1,36\sqrt D/\varepsilon\}},\qquad \widetilde T_\varepsilon=\frac Lm\max\{1,36\sqrt D/\varepsilon\}.\]


If \(N\varepsilon^2\ge18d_\theta\) and \(k\ge\widetilde T_\varepsilon\log(1+3\sqrt mR/\varepsilon)\), then

\[\sqrt{\frac{2d_\theta}{N}}+\sqrt m\left((1-m\widetilde\gamma_\varepsilon)^kR+2B(\widetilde\gamma_\varepsilon)\right)\le\varepsilon.\]


\textbf{Proof.} The new step and prefactor are the preceding definitions with accuracy \(\varepsilon/2\); simplifying their formulas gives the coefficient \(36\). The earlier bias budget at that half accuracy gives \(\sqrt mB\le\varepsilon/6\), so twice the bias costs at most \(\varepsilon/3\). The same reciprocal identity proves the transient budget \(\varepsilon/3\) at the displayed logarithmic threshold. The unchanged concentration budget is at most \(\varepsilon/3\). Add the three terms. This is a valid auxiliary budget for a doubled-bias estimate. The primary application proves a sharper exact-Gibbs initialization bound with a single bias and therefore uses the coefficient \(18\), rather than requiring this conservative coefficient \(36\).

\begin{leanbox}{Lean: \texttt{iplaTargetAccuracyStep}, \texttt{iplaTargetAccuracyPrefactor}, \texttt{iplaTargetAccuracyStep\_\allowbreak{}eq} and 3 more}
\leansource{\texttt{iplaTargetAccuracyStep} \quad WassersteinMixing/\allowbreak{}IPLARates.lean, lines 37–38}
\begin{leancode}
\leanline{0}{def\ iplaTargetAccuracyStep\ (L\ D\ ε\ :\ ℝ)\ :\ ℝ\ :=}
\leanline{2}{iplaAccuracyStep\ L\ D\ (ε\ /\ 2)}
\end{leancode}
\leansource{\texttt{iplaTargetAccuracyPrefactor} \quad WassersteinMixing/\allowbreak{}IPLARates.lean, lines 41–42}
\begin{leancode}
\leanline{0}{def\ iplaTargetAccuracyPrefactor\ (m\ L\ D\ ε\ :\ ℝ)\ :\ ℝ\ :=}
\leanline{2}{iplaAccuracyPrefactor\ m\ L\ D\ (ε\ /\ 2)}
\end{leancode}
\leansource{\texttt{iplaTargetAccuracyStep\_\allowbreak{}eq} \quad WassersteinMixing/\allowbreak{}IPLARates.lean, lines 176–182}
\begin{leancode}
\leanline{0}{theorem\ iplaTargetAccuracyStep\_eq\ (L\ D\ ε\ :\ ℝ)\ :}
\leanline{4}{iplaTargetAccuracyStep\ L\ D\ ε\ =}
\leanline{6}{1\ /\ (L\ *\ max\ 1\ (36\ *\ Real.sqrt\ D\ /\ ε))}
\end{leancode}
\leansource{\texttt{iplaTargetAccuracyPrefactor\_\allowbreak{}eq} \quad WassersteinMixing/\allowbreak{}IPLARates.lean, lines 185–191}
\begin{leancode}
\leanline{0}{theorem\ iplaTargetAccuracyPrefactor\_eq\ (m\ L\ D\ ε\ :\ ℝ)\ :}
\leanline{4}{iplaTargetAccuracyPrefactor\ m\ L\ D\ ε\ =}
\leanline{6}{(L\ /\ m)\ *\ max\ 1\ (36\ *\ Real.sqrt\ D\ /\ ε)}
\end{leancode}
\leansource{\texttt{ipla\_\allowbreak{}target\_\allowbreak{}transient\_\allowbreak{}budget} \quad WassersteinMixing/\allowbreak{}IPLARates.lean, lines 194–217}
\begin{leancode}
\leanline{0}{theorem\ ipla\_target\_transient\_budget\ (m\ L\ D\ R\ ε\ :\ ℝ)\ (hm\ :\ 0\ <\ m)}
\leanline{4}{(hmL\ :\ m\ ≤\ L)\ (hR\ :\ 0\ ≤\ R)\ (hε\ :\ 0\ <\ ε)\ (k\ :\ ℕ)}
\leanline{4}{(hk\ :\ iplaTargetAccuracyPrefactor\ m\ L\ D\ ε\ *}
\leanline{6}{Real.log\ (1\ +\ 3\ *\ Real.sqrt\ m\ *\ R\ /\ ε)\ ≤\ (k\ :\ ℝ))\ :}
\leanline{4}{Real.sqrt\ m\ *\ ((1\ -\ m\ *\ iplaTargetAccuracyStep\ L\ D\ ε)\ \textasciicircum{}\ k\ *\ R)\ ≤\ ε\ /\ 3}
\end{leancode}
\leansource{\texttt{ipla\_\allowbreak{}accuracy\_\allowbreak{}target\_\allowbreak{}arithmetic} \quad WassersteinMixing/\allowbreak{}IPLARates.lean, lines 220–233}
\begin{leancode}
\leanline{0}{theorem\ ipla\_accuracy\_target\_arithmetic\ (N\ dθ\ :\ ℕ)\ (hN\ :\ 0\ <\ N)\ (m\ L\ D\ R\ ε\ :\ ℝ)}
\leanline{4}{(hm\ :\ 0\ <\ m)\ (hmL\ :\ m\ ≤\ L)\ (hD\ :\ 0\ ≤\ D)\ (hR\ :\ 0\ ≤\ R)\ (hε\ :\ 0\ <\ ε)}
\leanline{4}{(hparticles\ :\ 18\ *\ (dθ\ :\ ℝ)\ ≤\ (N\ :\ ℝ)\ *\ ε\ \textasciicircum{}\ 2)\ (k\ :\ ℕ)}
\leanline{4}{(hk\ :\ iplaTargetAccuracyPrefactor\ m\ L\ D\ ε\ *}
\leanline{6}{Real.log\ (1\ +\ 3\ *\ Real.sqrt\ m\ *\ R\ /\ ε)\ ≤\ (k\ :\ ℝ))\ :}
\leanline{4}{Real.sqrt\ (2\ *\ (dθ\ :\ ℝ)\ /\ (N\ :\ ℝ))\ +}
\leanline{6}{Real.sqrt\ m\ *\ ((1\ -\ m\ *\ iplaTargetAccuracyStep\ L\ D\ ε)\ \textasciicircum{}\ k\ *\ R\ +}
\leanline{8}{2\ *\ iplaLinearBias\ m\ L\ (iplaTargetAccuracyStep\ L\ D\ ε)\ D)\ ≤\ ε}
\end{leancode}
\leanlabel{How the Lean proof follows the argument}
\texttt{ipla\allowbreak{}Target\allowbreak{}Accuracy\allowbreak{}Step} and \texttt{ipla\allowbreak{}Target\allowbreak{}Accuracy\allowbreak{}Prefactor} reuse the base definitions at half accuracy, while the two equality lemmas expose the exact coefficient. \texttt{ipla\_\allowbreak{}target\_\allowbreak{}transient\_\allowbreak{}budget} proves the logarithmic bound for that step without changing the target error budget. \texttt{ipla\_\allowbreak{}accuracy\_\allowbreak{}target\_\allowbreak{}arithmetic} combines it with the original concentration budget and twice the half-accuracy bias estimate.
\end{leanbox}

\section{\texorpdfstring{An exact regime where the discretization estimate improves Project 03}{An exact regime where the discretization estimate improves Project 03}}


\textbf{Claim.} For \(m>0\), \(L,D,\gamma\ge0\),

\[m\gamma\le\frac{121}{1600}\ \Longrightarrow\ B(\gamma)\le B_{03}(\gamma),\]
\[\quad m\gamma\le\frac{121}{6400}\ \Longrightarrow\ 2B(\gamma)\le B_{03}(\gamma).\]


\textbf{Proof.} Factor the new term as

\[B(\gamma)=\bigl(6\gamma\sqrt m\bigr)\left(\frac Lm\sqrt D\right).\]


The second factor is nonnegative and is shared with the Project 03 expression. Comparing the first factors reduces to \(6\gamma\sqrt m\le(33/20)\sqrt\gamma\); squaring the nonnegative sides and using the displayed bound on \(m\gamma\) proves it. Replacing \(6\) by \(12\) gives the second threshold and the doubled-bias result. No division by \(\gamma\) or \(D\) is required, so their zero cases are included.




The primary application attains a single bias with either discretized-invariant or exact-Gibbs initialization, so its sharp complete bound uses the first threshold. The doubled statement records a valid conservative alternative. Thus the improved discretization order is linear in the step, within the explicitly stated regularity and step regime, rather than a claim that every Project 03 hypothesis or every allowed step is covered.

\begin{leanbox}{Lean: \texttt{ipla\_\allowbreak{}bias\_\allowbreak{}coefficient\_\allowbreak{}le\_\allowbreak{}old}, \texttt{ipla\_\allowbreak{}doubled\_\allowbreak{}bias\_\allowbreak{}coefficient\_\allowbreak{}le\_\allowbreak{}old}, \texttt{ipla\_\allowbreak{}linear\_\allowbreak{}bias\_\allowbreak{}factorization} and 2 more}
\leansource{\texttt{ipla\_\allowbreak{}bias\_\allowbreak{}coefficient\_\allowbreak{}le\_\allowbreak{}old} \quad WassersteinMixing/\allowbreak{}IPLARates.lean, lines 236–242}
\begin{leancode}
\leanline{0}{theorem\ ipla\_bias\_coefficient\_le\_old\ (m\ γ\ :\ ℝ)\ (hm\ :\ 0\ <\ m)\ (hγ\ :\ 0\ ≤\ γ)}
\leanline{4}{(hsmall\ :\ m\ *\ γ\ ≤\ (121\ /\ 1600\ :\ ℝ))\ :}
\leanline{4}{6\ *\ γ\ *\ Real.sqrt\ m\ ≤\ (33\ /\ 20\ :\ ℝ)\ *\ Real.sqrt\ γ}
\end{leancode}
\leansource{\texttt{ipla\_\allowbreak{}doubled\_\allowbreak{}bias\_\allowbreak{}coefficient\_\allowbreak{}le\_\allowbreak{}old} \quad WassersteinMixing/\allowbreak{}IPLARates.lean, lines 245–251}
\begin{leancode}
\leanline{0}{theorem\ ipla\_doubled\_bias\_coefficient\_le\_old\ (m\ γ\ :\ ℝ)\ (hm\ :\ 0\ <\ m)\ (hγ\ :\ 0\ ≤\ γ)}
\leanline{4}{(hsmall\ :\ m\ *\ γ\ ≤\ (121\ /\ 6400\ :\ ℝ))\ :}
\leanline{4}{12\ *\ γ\ *\ Real.sqrt\ m\ ≤\ (33\ /\ 20\ :\ ℝ)\ *\ Real.sqrt\ γ}
\end{leancode}
\leansource{\texttt{ipla\_\allowbreak{}linear\_\allowbreak{}bias\_\allowbreak{}factorization} \quad WassersteinMixing/\allowbreak{}IPLARates.lean, lines 254–269}
\begin{leancode}
\leanline{0}{theorem\ ipla\_linear\_bias\_factorization\ (m\ L\ γ\ D\ :\ ℝ)\ (hm\ :\ 0\ <\ m)}
\leanline{4}{(hL\ :\ 0\ ≤\ L)\ (hD\ :\ 0\ ≤\ D)\ :}
\leanline{4}{iplaLinearBias\ m\ L\ γ\ D\ =}
\leanline{6}{(6\ *\ γ\ *\ Real.sqrt\ m)\ *\ ((L\ /\ m)\ *\ Real.sqrt\ D)}
\end{leancode}
\leansource{\texttt{ipla\_\allowbreak{}linear\_\allowbreak{}bias\_\allowbreak{}le\_\allowbreak{}project03} \quad WassersteinMixing/\allowbreak{}IPLARates.lean, lines 273–282}
\begin{leancode}
\leanline{0}{theorem\ ipla\_linear\_bias\_le\_project03\ (m\ L\ γ\ D\ :\ ℝ)\ (hm\ :\ 0\ <\ m)}
\leanline{4}{(hL\ :\ 0\ ≤\ L)\ (hγ\ :\ 0\ ≤\ γ)\ (hD\ :\ 0\ ≤\ D)}
\leanline{4}{(hsmall\ :\ m\ *\ γ\ ≤\ (121\ /\ 1600\ :\ ℝ))\ :}
\leanline{4}{iplaLinearBias\ m\ L\ γ\ D\ ≤\ iplaProject03Bias\ m\ L\ γ\ D}
\end{leancode}
\leansource{\texttt{ipla\_\allowbreak{}doubled\_\allowbreak{}bias\_\allowbreak{}le\_\allowbreak{}project03} \quad WassersteinMixing/\allowbreak{}IPLARates.lean, lines 287–296}
\begin{leancode}
\leanline{0}{theorem\ ipla\_doubled\_bias\_le\_project03\ (m\ L\ γ\ D\ :\ ℝ)\ (hm\ :\ 0\ <\ m)}
\leanline{4}{(hL\ :\ 0\ ≤\ L)\ (hγ\ :\ 0\ ≤\ γ)\ (hD\ :\ 0\ ≤\ D)}
\leanline{4}{(hsmall\ :\ m\ *\ γ\ ≤\ (121\ /\ 6400\ :\ ℝ))\ :}
\leanline{4}{2\ *\ iplaLinearBias\ m\ L\ γ\ D\ ≤\ iplaProject03Bias\ m\ L\ γ\ D}
\end{leancode}
\leanlabel{How the Lean proof follows the argument}
\texttt{ipla\_\allowbreak{}bias\_\allowbreak{}coefficient\_\allowbreak{}le\_\allowbreak{}old} and \texttt{ipla\_\allowbreak{}doubled\_\allowbreak{}bias\_\allowbreak{}coefficient\_\allowbreak{}le\_\allowbreak{}old} prove the scalar comparisons by nonnegative-square inequalities. \texttt{ipla\_\allowbreak{}linear\_\allowbreak{}bias\_\allowbreak{}factorization} derives the common-factor identity from the curvature square-root scaling. The final two lemmas multiply the appropriate scalar inequality by the nonnegative common factor and rewrite the exact Project 03 expression.
\end{leanbox}

\end{document}
